%==============================================================================
\chapter{Results and Discussion}
\label{ch:results}
%==============================================================================

This chapter reports what the running system produces and what was verified,
and is explicit about what has not yet been measured. The claims below are
about the realized system's observed behaviour; no user-study or
learning-outcome figures are reported, because none were collected.

%------------------------------------------------------------------------------
\section{Results and Analysis}
\label{sec:results}
%------------------------------------------------------------------------------

\subsection{End-to-End Generation Against Live Models}

With the real Gemini providers configured, the pipeline produced a complete
lesson for the topic \emph{Photosynthesis}: six narration segments totalling
approximately forty-three seconds of audio, a well-formed SVG whose parts
carry semantic identifiers such as \texttt{sun}, \texttt{leaf}, and
\texttt{chloroplast}, narration whose references all passed the
deterministic validator, and real synthesized speech for each segment. This
single run exercises every node of the graph --- planner, SVG generator,
critique/refine, narration generator, validator, TTS renderer, and assembler
--- against live models. Figure~\ref{fig:diagram} shows the generated
diagram as rendered by the app.

\begin{figure}[htbp]
\centering
\shot[width=0.85\textwidth]{diagram.png}
\caption{The diagram generated for \emph{Photosynthesis}, as shown in the
app's full-screen view. Each labelled part --- the sun, the leaf, the
chloroplast, the water and carbon-dioxide inputs, and the glucose and oxygen
outputs --- is a separate SVG element with a stable identifier, which is
exactly what lets the narration highlight it during playback.}
\label{fig:diagram}
\end{figure}

The saved library shown in Figure~\ref{fig:screens} holds further lessons
generated for unrelated topics --- the Pythagorean theorem and the
Transformer architecture --- demonstrating that the pipeline is not tuned to
a single subject: the same prompts, graph, and validator produced consistent
lessons across a natural-science topic, a mathematical topic, and a
computer-science topic.

\begin{figure}[htbp]
\centering
\begin{tabular}{ccc}
\shot[width=0.30\textwidth]{home.png} &
\shot[width=0.30\textwidth]{player.png} &
\shot[width=0.30\textwidth]{library.png}\\
(a) Topic entry & (b) Player & (c) Library\\
\end{tabular}
\caption{The Android client. (a)~The home screen where a topic is typed.
(b)~The player mid-lesson: the current segment's diagram parts are
highlighted while its audio plays, with the caption below. (c)~The on-device
library of saved lessons, each replayable fully offline.}
\label{fig:screens}
\end{figure}

\subsection{The Validation Gate Works in Both Directions}

The property the design promises is that the gate both catches broken
lessons and does not block valid ones. Both directions were verified by the
graph tests:

\begin{itemize}
  \item a lesson containing a deliberately broken part reference is
        \emph{rejected} by the validator and routed into the repair loop,
        and when the repair succeeds the corrected lesson passes and
        proceeds to rendering; and
  \item when repair cannot fix the problem within the retry budget $R$, the
        job \emph{fails cleanly} --- no partial or inconsistent lesson is
        ever cached.
\end{itemize}

Combined with the live run above (where valid generations passed the gate
unimpeded), this confirms the system-wide invariant: every lesson that
reaches the cache satisfies ``all referenced identifiers exist.''

\subsection{Synchronization Behaviour}

Because each segment's duration is computed from the synthesized audio's
byte count and sample rate, and the on-device highlight is driven by the
audio player's current-item index, the highlight transition coincides with
the clip transition by construction. Playback, pause, resume,
seek-by-segment, and replay all preserve synchronization, since the only
synchronization state is the current clip index. The same behaviour holds
for offline replays, where all assets are local files.

\subsection{Reproducibility and Automation Results}

Table~\ref{tab:verified} summarizes the verification of the engineering
infrastructure.

\begin{table}[htbp]
\caption{Verified engineering results}
\label{tab:verified}
\centering
\small
\renewcommand{\arraystretch}{1.3}
\begin{tabular}{@{}p{0.34\textwidth}p{0.58\textwidth}@{}}
\toprule
\textbf{Aspect} & \textbf{Verified result}\\
\midrule
Backend test suite & 48 hermetic tests pass with \emph{zero} live API calls
and no keys present\\
Container & Docker image builds; container serves \texttt{/healthz} (200);
mock-provider generation returns the expected 202 asynchronous response with
no key\\
Client build & Debug APK assembles reproducibly in CI; release workflow
produces a signed APK\\
Deployment & Backend image published to a registry and deployed to a cloud
host behind auto-HTTPS by the deploy workflow\\
Offline path & Save-offline persists all assets locally; the player loads
saved lessons from local file URIs through the same code path as fresh
lessons\\
Model-id resilience & Live model listing tool confirms pinned identifiers;
changing models is a settings change only\\
\bottomrule
\end{tabular}
\end{table}

\subsection{The Realized System at a Glance}

\begin{table}[htbp]
\caption{Realized system summary}
\label{tab:summary}
\centering
\small
\renewcommand{\arraystretch}{1.3}
\begin{tabular}{@{}ll@{}}
\toprule
\textbf{Aspect} & \textbf{Realization}\\
\midrule
Backend & FastAPI, $\approx$2.0k LoC Python\\
Client & Android/Compose, $\approx$3.8k LoC Kotlin\\
Generation & LangGraph: plan $\to$ draw $\to$ critique/refine $\to$ narrate
$\to$ validate $\to$ repair\\
Gate & Deterministic non-LLM validator\\
Models & Text + speech via swappable adapters (Gemini Flash class)\\
Caching & Stateless, topic-keyed, coalesced jobs\\
Offline & Room + file store; network-free replay\\
Tests & 48 hermetic, key-free backend tests\\
Delivery & Container + CI/CD + auto-HTTPS deployment\\
\bottomrule
\end{tabular}
\end{table}

%------------------------------------------------------------------------------
\section{Discussion}
\label{sec:discussion}
%------------------------------------------------------------------------------

\subsection{Analysis of the Design Choices}

\textbf{The gate, not the model, carries the guarantee.} The central result
of this work is architectural: the correctness-critical property (narration
references resolve) is enforced by a component that cannot hallucinate. The
LLM-facing mechanisms --- planning first, narrating against the finished
drawing, structured output, low temperature --- \emph{reduce the frequency}
of violations, and the repair loop \emph{recovers} from most of the rest,
but the deterministic validator is what converts ``usually consistent'' into
``never served inconsistent.'' This separation of concerns (probabilistic
quality mechanisms in front of a deterministic correctness mechanism) is
transferable to any system whose LLM outputs must satisfy checkable
invariants.

\textbf{Bounded loops keep the worst case owned.} Both self-correction loops
carry explicit budgets, so the maximum number of model calls per lesson is
known in advance, and the failure mode of an unfixable topic is a clean,
retryable error rather than an unbounded spend. This is a deliberate
divergence from open-ended agent designs in the literature.

\textbf{The cache changes the economics.} Statelessness plus topic-keyed
caching means the marginal cost of a popular topic tends to zero and the
response time collapses to a lookup. Request coalescing caps the cost spike
of a suddenly popular topic at one generation.

\subsection{Limitations}

Several limitations are worth naming plainly.

\emph{Diagram quality has a ceiling.} A language model drawing SVG is good
at clear, well-known concepts and weaker at dense or unusual ones. The
critique-and-refine loop raises the floor but does not guarantee a beautiful
result: the deterministic gate guarantees \emph{consistency}, not aesthetic
quality --- a valid lesson can still be a plain one.

\emph{Synchronization is at the segment level.} A group of parts is lit for
a whole sentence, not each part at its spoken word. This was a deliberate
trade for determinism and offline play; word-level highlighting would be
more precise but would reintroduce timing data and fragility.

\emph{Cold topics are slow and cost money.} The first request for a new
topic pays for several model calls and speech synthesis and takes seconds.
The cache hides this for repeated topics and the bounded loops cap the worst
case, but the cold path is inherently the slow path.

\emph{Generated markup must be treated as untrusted.} Rendering
model-produced SVG in a WebView requires care; the WebView is restricted to
the app's own highlight script and content-supplied script is not executed,
but this remains an area demanding ongoing attention.

\emph{Evaluation breadth.} The claims here are about verified system
behaviour, not measured learning gains. No controlled user study, no large
quantitative study of diagram quality across many topics, and no on-device
latency or battery measurements across a fleet of devices have been
performed; these are flagged as future work rather than claimed as results.
